%% ══════════════════════════════════════════════════════════════
%%  CHAPTER 3 — AGENTIC ARCHITECTURE DEEP DIVE
%% ══════════════════════════════════════════════════════════════
\chapter{Agentic Architecture}
\label{ch:agents}

\section{Overview of the Multi-Agent System}
\label{sec:agents:overview}

\autoscholar employs a \emph{prompt-based multi-agent architecture} where a single \gls{llm} instance (\claude) assumes different roles throughout the research pipeline.
Unlike frameworks where each agent is a separate model instance with its own context, \autoscholar's agents are conceptual personas defined in the system prompt (\file{server/services/claude/systemPrompt.ts}, 1,798 lines).
Claude switches between these personas based on the current pipeline stage, guided by explicit instructions in the system prompt and reinforced by \governor directives.

The architecture comprises seven distinct agent roles:

\begin{table}[H]
\centering
\caption{Agent roles in the \autoscholar system.}
\label{tab:agent-roles}
\begin{tabularx}{\textwidth}{llX}
\toprule
\textbf{Agent} & \textbf{Named After} & \textbf{Role} \\
\midrule
\euler & Leonhard Euler & Global orchestrator, academic writer, \LaTeX{} compiler \\
\gauss & Carl Friedrich Gauss & Literature gap analysis and research question formulation \\
\turing & Alan Turing & Data collection from APIs, databases, and web sources \\
\newton & Isaac Newton & Dataset preparation, cleaning, feature engineering \\
\fisher & Ronald Fisher & Empirical analysis, econometric modeling, figure generation \\
\governor & --- & Rule-based pipeline supervisor (no LLM) \\
\popper & Karl Popper & AI-powered peer review and quality validation \\
\bottomrule
\end{tabularx}
\end{table}

\section{Euler --- The Global Orchestrator}
\label{sec:agents:euler}

\subsection{Role and Responsibilities}

\euler is the master agent that coordinates the entire research pipeline.
It is the ``default mode'' of Claude throughout the session and has three primary responsibilities:

\begin{enumerate}
  \item \textbf{Pipeline orchestration.} \euler decides when to transition between pipeline phases (literature review $\to$ data collection $\to$ analysis $\to$ writing $\to$ compilation).
  \item \textbf{Academic writing.} \euler authors all paper sections in \LaTeX{}, following the Journal of Finance writing style directive (282 lines of style instructions).
  \item \textbf{LaTeX compilation.} \euler manages the paper compilation process, including section assembly, BibTeX generation, and iterative error repair.
\end{enumerate}

\subsection{Decision-Making Process}

\euler's decision-making is governed by the system prompt's 7-step research pipeline:

\begin{algorithm}[H]
\caption{Euler's Research Pipeline Orchestration}
\label{alg:euler-pipeline}
\SetAlgoLined
\KwIn{Research prompt $p$, available tools $\mathcal{T}$}
\KwOut{Compiled PDF paper, project ZIP}

\textbf{Step 1: Gap Analysis} (invoke \gauss)\;
Search academic databases for existing literature\;
Identify research gaps and formulate hypotheses\;

\textbf{Step 2: Data Collection} (invoke \turing)\;
Query financial APIs, macroeconomic databases, web sources\;
Download datasets to the Python environment\;

\textbf{Step 3: Data Preparation} (invoke \newton)\;
Clean, merge, and transform raw datasets\;
Construct derived variables and panel structures\;

\textbf{Step 4: Empirical Analysis} (invoke \fisher)\;
Estimate models, run robustness checks\;
Generate publication-quality figures and tables\;

\textbf{Step 5: Paper Writing} (invoke \euler writing mode)\;
Write each section using \texttt{write\_paper\_section} tool\;
Target 60--90 pages of content\;

\textbf{Step 6: Compilation} (invoke \euler compilation mode)\;
Compile \LaTeX{} via \texttt{compile\_research\_paper}\;
If errors, enter repair loop (read $\to$ patch $\to$ recompile)\;

\textbf{Step 7: Packaging}\;
Create reproducible project ZIP via \texttt{compile\_research\_paper}\;
\end{algorithm}

\subsection{State Transitions}

\euler does not maintain explicit state---its ``state'' is the conversation history and the accumulated tool results within the 1M-token context window.
State transitions are implicit: when \euler finishes analyzing literature, it naturally transitions to data collection by invoking \turing-associated tools.

The system prompt provides explicit ordering constraints:
\begin{quote}
\itshape
``The 7-step pipeline is MANDATORY. You cannot write a paper section before collecting data, and you cannot collect data before identifying what data you need through gap analysis.''
\end{quote}

\subsection{Writing Mode}

When \euler enters writing mode, it follows the Journal of Finance style directive (\file{governor/jfWritingStyle.ts}, 282 lines), which includes:

\begin{itemize}
  \item Prose identity guidance from editors Campbell Harvey, John Cochrane, Andrei Shleifer, Lubos Pastor, and Craig Holden.
  \item Official JF formatting rules (100-word abstract limit, Roman numeral tables, unnumbered introduction).
  \item A banned language registry with 100+ banned words/phrases and tier-based restrictions.
  \item Section-by-section writing conventions.
  \item Advanced prose craft techniques (emphatic stress, topic-comment structure).
  \item AI detection countermeasures (lexical diversity, syntactic variation).
  \item An 8-pass verification checklist.
\end{itemize}

\subsection{Tool Access}

\euler has access to all 213 tools, but primarily uses:

\begin{itemize}
  \item \texttt{write\_paper\_section} --- Write individual paper sections in \LaTeX{}.
  \item \texttt{compile\_research\_paper} --- Assemble sections and compile to PDF.
  \item \texttt{read\_latex\_source} / \texttt{inspect\_latex\_errors} / \texttt{patch\_latex} / \texttt{recompile\_latex} --- The 4-step \LaTeX{} repair loop.
  \item \texttt{generate\_research\_paper} --- Direct single-call paper generation (legacy path).
\end{itemize}

\section{Gauss --- The Gap Finder}
\label{sec:agents:gauss}

\subsection{Role and Responsibilities}

\gauss is the literature discovery and gap analysis agent, responsible for identifying what research has been done and---critically---what has \emph{not} been done.
It is implemented as both a conceptual role in the system prompt and a standalone service (\file{literatureGapFinderService.ts}, 811 lines).

\subsection{The 7-Step Gap Discovery Pipeline}

\gauss implements a sophisticated multi-stage pipeline:

\begin{description}
  \item[Step 1: Search Strategy Planning.] Claude generates 7 diverse search queries from different angles for the target topic. This ensures broad coverage rather than narrow keyword matching.

  \item[Step 2: Multi-Source Parallel Search.] Queries are dispatched simultaneously to 4 academic databases:
  \begin{itemize}
    \item \textbf{Semantic Scholar} --- AI-powered academic search with citation data, retry with exponential backoff on rate limits.
    \item \textbf{OpenAlex} --- Open bibliographic database with 200M+ works.
    \item \textbf{arXiv} --- Preprint server for cutting-edge research.
    \item \textbf{CrossRef} --- DOI registration agency with 130M+ records.
  \end{itemize}
  Results are standardized into \texttt{AcademicPaper} objects with auto-generated BibTeX.

  \item[Step 3: Citation Network Expansion.] For high-impact papers ($>$50 citations), \gauss retrieves their references and citing papers via OpenAlex and CrossRef, expanding the literature landscape.

  \item[Step 4: Scientific Structure Extraction.] Claude analyzes each paper to extract:
  \begin{itemize}
    \item Methods used (econometric, statistical, ML)
    \item Datasets employed
    \item Key findings
    \item Stated limitations
    \item Temporal coverage
  \end{itemize}

  \item[Step 5: Literature Synthesis.] Claude clusters papers into thematic groups and identifies:
  \begin{itemize}
    \item Consensus findings
    \item Contradictions between papers
    \item Methodological trends
    \item Temporal evolution of the field
  \end{itemize}

  \item[Step 6: Gap Detection.] Six specialized gap detection engines run in parallel:
  \begin{enumerate}
    \item \textbf{Method gaps} --- Methods used in adjacent fields but not applied to this topic.
    \item \textbf{Dataset gaps} --- Available datasets that have not been exploited.
    \item \textbf{Temporal gaps} --- Time periods or recent data not yet studied.
    \item \textbf{Contradiction gaps} --- Conflicting findings that need resolution.
    \item \textbf{Limitation gaps} --- Explicitly stated limitations in existing papers.
    \item \textbf{Cross-disciplinary gaps} --- Insights from other fields applicable here.
  \end{enumerate}

  \item[Step 7: Gap Validation and Scoring.] Each identified gap is scored for novelty, feasibility, and impact, producing a ranked list of research opportunities.
\end{description}

\subsection{Robust JSON Parsing}

A critical implementation detail is the 6-strategy JSON parsing system in \func{parseJsonResponse}. Since LLM outputs frequently contain malformed JSON (trailing commas, unescaped quotes, markdown code fences), the parser attempts:

\begin{enumerate}
  \item Direct \texttt{JSON.parse}.
  \item Strip markdown code fences and re-parse.
  \item Extract JSON from mixed text using regex.
  \item Fix common issues (trailing commas, single quotes).
  \item Parse partial JSON up to the last valid closing bracket.
  \item Return a fallback structure with the raw text.
\end{enumerate}

\subsection{External Tools Used}

\begin{itemize}
  \item \texttt{academic\_search\_multi} --- Parallel search across all 5 academic databases.
  \item \texttt{literature\_gap\_finder} --- The full 7-step pipeline.
  \item \texttt{semantic\_scholar\_search}, \texttt{arxiv\_search}, \texttt{crossref\_lookup}, \texttt{openalex\_search} --- Individual source queries.
\end{itemize}

\subsection{Failure Modes}

\begin{itemize}
  \item \textbf{Rate limiting:} Semantic Scholar frequently returns 429s; the system uses exponential backoff (up to 3 retries with 1s, 2s, 4s delays).
  \item \textbf{Sparse results:} For niche topics, some databases may return zero results. The multi-source approach ensures at least some results are found.
  \item \textbf{Malformed LLM output:} The 6-strategy JSON parser handles most cases, but extreme malformation can produce degraded gap analysis.
\end{itemize}

\section{Turing --- The Data Collector}
\label{sec:agents:turing}

\subsection{Role and Responsibilities}

\turing is responsible for acquiring all raw data needed for the research paper.
This encompasses financial market data, macroeconomic indicators, academic datasets, and web-sourced information.

\subsection{Data Source Hierarchy}

The system prompt defines a strict data source hierarchy with fallback chains:

\begin{enumerate}
  \item \textbf{Primary: FMP API} (230+ endpoints) --- Company fundamentals, market quotes, historical prices, technical indicators, SEC filings, insider trading, ESG scores, earnings data, analyst estimates, and more.
  \item \textbf{Secondary: EODHD} --- Extended historical data back to IPO date, intraday bars, US options chains with Greeks.
  \item \textbf{Tertiary: FRED} --- 840,000+ macroeconomic time series with unit transformations.
  \item \textbf{Macro: World Bank / IMF / OECD} --- Free, keyless access to 16,000+ development indicators, WEO data, and 1,475+ OECD dataflows.
  \item \textbf{Web: Tavily / Firecrawl / Exa} --- When structured APIs lack the required data, web search and extraction fill the gap.
  \item \textbf{Prediction Markets: Polymarket} --- Implied probabilities for event-driven research.
\end{enumerate}

\subsection{Data Collection Methods}

\turing uses three classes of data collection tools:

\begin{description}
  \item[Direct API tools.] 170+ tools that call specific API endpoints (e.g., \texttt{get\_income\_statement}, \texttt{get\_wb\_indicator\_data}, \texttt{get\_fred\_series}). Each returns structured JSON.

  \item[Python data clients.] 5 pre-loaded Python clients (\file{fmp\_client.py}, \file{fred\_client.py}, \file{eodhd\_client.py}, \file{macro\_clients.py}) enable DataFrame-native data acquisition directly in analysis scripts.

  \item[Web data extraction.] The \gls{swdes} (\file{webDatasetService.ts}, 1,431 lines) provides autonomous large-scale web data collection:
  \begin{enumerate}
    \item Claude Haiku plans search queries and extraction schemas.
    \item Parallel multi-source search across Firecrawl, Tavily, and Exa.
    \item 6-strategy extraction pipeline (schema-based, portal crawling, batch scraping, download links, fallback).
    \item Schema enforcement and deduplication via Claude Haiku.
    \item Statistical validation and completeness checks.
    \item CSV/JSON output with metadata sidecar.
  \end{enumerate}
  Capable of collecting up to 100,000 rows from 25+ sources in a single call.
\end{description}

\subsection{Data Validation}

\turing validates data at collection time:
\begin{itemize}
  \item API response schema validation.
  \item Null/NaN sanitization (especially for EODHD options data).
  \item Ticker normalization (e.g., ``AAPL'' $\to$ ``AAPL.US'' for EODHD).
  \item Date format handling (Unix timestamps, ISO strings, datetime objects).
  \item Result size limiting to prevent context overflow (e.g., historical prices capped at 504 points).
\end{itemize}

\subsection{Failure Modes}

\begin{itemize}
  \item \textbf{API key exhaustion:} FMP has daily request limits; the system has no automatic quota management.
  \item \textbf{Data gaps:} Some tickers or time periods may have missing data; the system relies on Claude to detect and handle gaps.
  \item \textbf{Timeout:} External API calls use 15--60s timeouts; slow services may cause tool execution failures.
\end{itemize}

\section{Newton --- The Data Preparer}
\label{sec:agents:newton}

\subsection{Role and Responsibilities}

\newton transforms raw data into analysis-ready datasets.
This involves cleaning, merging, feature construction, and quality validation.

\subsection{Data Preparation Pipeline}

\newton operates primarily through the \texttt{execute\_custom\_python\_analysis} tool, running Python code that:

\begin{enumerate}
  \item \textbf{Cleans raw data.} Handles missing values (forward fill, interpolation, or explicit NaN treatment based on data type), removes duplicates, and corrects data types.

  \item \textbf{Merges datasets.} Aligns time series by date, merges cross-sectional data by identifiers, and handles frequency mismatches (e.g., daily prices with quarterly fundamentals).

  \item \textbf{Constructs derived variables.} Calculates returns (simple, log, excess), volatility measures, momentum indicators, financial ratios, and interaction terms.

  \item \textbf{Builds panel structures.} Constructs entity $\times$ time panels with proper indexing for fixed-effects regression.

  \item \textbf{Validates quality.} Checks for:
  \begin{itemize}
    \item Stationarity (ADF/KPSS tests)
    \item Multicollinearity (VIF)
    \item Outliers (winsorization at 1\%/99\%)
    \item Missing data patterns (MCAR/MAR/MNAR)
    \item Sample size adequacy
  \end{itemize}
\end{enumerate}

\subsection{Automatic Error Patching}

The \file{customPython.ts} module applies several automatic fixes to \newton's code before execution:

\begin{itemize}
  \item \textbf{Deprecated pandas frequencies:} Replaces \texttt{'M'} with \texttt{'ME'}, \texttt{'Q'} with \texttt{'QE'}, etc.
  \item \textbf{NumPy deprecations:} Replaces \texttt{np.trapz} with \texttt{np.trapezoid}, removes \texttt{np.math} references.
  \item \textbf{Scikit-learn API changes:} Renames \texttt{n\_iter} to \texttt{max\_iter} for t-SNE.
  \item \textbf{Matplotlib conflicts:} Patches \texttt{tight\_layout} to handle colorbar conflicts.
  \item \textbf{Data fabrication detection:} Flags scripts that generate random data without API calls.
\end{itemize}

\subsection{Failure Modes}

\begin{itemize}
  \item \textbf{Memory errors:} Large datasets can exceed the Python subprocess memory limit.
  \item \textbf{Key errors:} Missing columns due to API response schema changes.
  \item \textbf{Type mismatches:} Mixed types in columns causing downstream computation failures.
\end{itemize}

\section{Fisher --- The Analysis Engine}
\label{sec:agents:fisher}

\subsection{Role and Responsibilities}

\fisher conducts all empirical analysis: econometric estimation, statistical inference, machine learning modeling, and figure generation.
It is the most computationally intensive agent, heavily relying on the Python execution environment.

\subsection{Methodology Tracks}

The system prompt defines four methodology tracks, each with specific model inventories:

\subsubsection{Track 1: Econometrics}

\begin{itemize}
  \item OLS, WLS, GLS, 2SLS, GMM, panel fixed/random effects
  \item Time series: ARIMA, VAR, VECM, GARCH(p,q), DCC-GARCH
  \item Causal inference: difference-in-differences, regression discontinuity, instrumental variables, propensity score matching
  \item High-frequency: realized volatility, Amihud illiquidity, Kyle's $\lambda$
\end{itemize}

\subsubsection{Track 2: Statistics}

\begin{itemize}
  \item Hypothesis testing: t-tests, F-tests, $\chi^2$ tests, Kolmogorov-Smirnov, Mann-Whitney
  \item Non-parametric: kernel density estimation, bootstrap confidence intervals
  \item Bayesian: MCMC, variational inference, Bayesian model averaging
  \item Distribution: quantile regression, extreme value theory, copulas
\end{itemize}

\subsubsection{Track 3: Machine Learning}

\begin{itemize}
  \item Supervised: random forests, gradient boosting (XGBoost, LightGBM), SVMs, elastic net
  \item Unsupervised: K-means, hierarchical clustering, PCA, t-SNE, UMAP
  \item Validation: cross-validation, grid/random search, learning curves
  \item Interpretability: SHAP values, partial dependence plots, feature importance
\end{itemize}

\subsubsection{Track 4: AI/Deep Learning}

\begin{itemize}
  \item Architectures: LSTM, GRU, Transformer, attention mechanisms
  \item NLP: sentiment analysis, named entity recognition, topic modeling
  \item Transfer learning: pre-trained embeddings, fine-tuning
  \item Evaluation: ablation studies, attention visualization
\end{itemize}

\subsection{Pre-Built Analysis Services}

In addition to custom Python execution, \fisher has access to 7 dedicated analysis services:

\begin{table}[H]
\centering
\caption{Pre-built quantitative analysis services.}
\label{tab:analysis-services}
\begin{tabularx}{\textwidth}{lXr}
\toprule
\textbf{Service} & \textbf{Capabilities} & \textbf{Lines} \\
\midrule
GARCH & GARCH(p,q) estimation, volatility forecasting, persistence metrics & 506 \\
VaR & Historical, parametric, Monte Carlo, and CVaR/Expected Shortfall & 555 \\
Monte Carlo & Investment projections with configurable simulations and horizons & 453 \\
Portfolio Optimizer & Markowitz MVO, efficient frontier, max Sharpe, min volatility & 672 \\
Risk Metrics & 20+ metrics: Sharpe, Sortino, Calmar, beta, alpha, drawdown & 795 \\
Options Pricing & Black-Scholes, Greeks, implied volatility (Newton-Raphson) & 589 \\
Volatility Surface & 3D IV surface construction with cubic interpolation & 520 \\
\bottomrule
\end{tabularx}
\end{table}

\subsection{Figure Generation}

\fisher generates publication-quality figures following the academic figure style directive (\file{governor/academicFigureStyle.ts}, 429 lines):

\begin{itemize}
  \item \textbf{Design philosophy:} Monochrome-first, print-first, Tufte data-ink ratio maximization.
  \item \textbf{Typography:} 9pt serif axis labels, 0.5pt spines, inward ticks.
  \item \textbf{Color policy:} Black $>$ grayscale $>$ Paul Tol/Okabe-Ito palettes. Rainbow/jet palettes are explicitly prohibited.
  \item \textbf{Chart conventions:} 16 chart types with specific formatting rules (time-series, event-study CAR, coefficient plots, bar charts, heatmaps, impulse response, RD plots, DiD, fan charts, survival curves).
  \item \textbf{Export:} 300 DPI PNG for inclusion, PDF with TrueType fonts for publication.
  \item \textbf{Watermark:} AutoScholar watermark added to all figures.
\end{itemize}

The matplotlib configuration (\file{figures/matplotlib\_config.py}, 79 lines) sets serif fonts, an 8-color muted palette, inward ticks, minimal chrome, and print-safe grid styling.
Blank figure detection (\file{figures/figure\_capture.py}, 104 lines) uses two-stage analysis: axes artist inspection followed by pixel-level variance measurement.

\subsection{Failure Modes}

\begin{itemize}
  \item \textbf{Convergence failures:} GARCH and maximum likelihood models may fail to converge with insufficient or degenerate data.
  \item \textbf{Singular matrices:} Covariance matrices with perfectly correlated or constant variables.
  \item \textbf{Memory errors:} Large Monte Carlo simulations or high-dimensional panel datasets.
  \item \textbf{Module import failures:} Missing Python packages (mitigated by the \texttt{install\_python\_package} tool).
\end{itemize}

\section{Governor --- The Pipeline Supervisor}
\label{sec:agents:governor}

\subsection{Architecture}

The \governor is unique among \autoscholar's agents: it is a \textbf{pure rule-based system} with zero LLM calls and zero database access.
Implemented in \file{governor/governorEngine.ts} (472 lines), it observes tool results in real-time and issues corrective actions when problems are detected.

\subsection{State Model}

The \governor maintains a typed state object (\file{governor/types.ts}, 164 lines):

\begin{lstlisting}[style=typescript, caption=Governor state model (simplified).]
interface AppState {
  runId: string;
  stage: PipelineStage;  // gauss | turing | newton | fisher | euler_write | euler_compile | popper | complete
  profile: string;       // active prompt profile
  pythonErrors: number;
  latexErrors: number;
  dataIssues: number;
  retryBudgets: {
    python: number;   // default: 3
    latex: number;     // default: 3
    data: number;      // default: 2
  };
  confidence: number;   // 0.0 - 1.0, starts at 1.0
  figureCount: number;
  sectionWordCounts: Map<string, number>;
  warnings: string[];
  publicationVerdict: PublicationVerdict | null;
}
\end{lstlisting}

\subsection{Decision Rules}

The \governor implements 7 decision rules, evaluated in priority order:

\begin{algorithm}[H]
\caption{Governor Decision Engine}
\label{alg:governor}
\SetAlgoLined
\KwIn{Batch of tool results $B$, current state $S$}
\KwOut{Governor evaluation (actions, updated state)}

\ForEach{tool result $r \in B$}{
  \textbf{Rule 1:} Detect pipeline stage from tool name, update $S.\text{stage}$\;
  \textbf{Rule 2:} If Python error detected:
    \begin{itemize}
      \item Parse error pattern (ModuleNotFoundError, MemoryError, KeyError, convergence, singular matrix)
      \item Generate pattern-specific recovery advice
      \item Decrement $S.\text{retryBudgets.python}$; if exhausted, emit \texttt{force\_finalize}
      \item Degrade $S.\text{confidence}$ by 0.15
      \item Switch to \texttt{debug\_recovery} profile
    \end{itemize}
  \textbf{Rule 3:} If data quality issue (empty dataset, missing columns, NaN ratio):
    \begin{itemize}
      \item Decrement $S.\text{retryBudgets.data}$
      \item Switch to \texttt{conservative} profile
    \end{itemize}
  \textbf{Rule 4:} If \LaTeX{} compilation failure:
    \begin{itemize}
      \item Decrement $S.\text{retryBudgets.latex}$
      \item Switch to \texttt{latex\_fallback} profile (restricts to basic \LaTeX{})
    \end{itemize}
  \textbf{Rule 5:} Track figure count from Python results\;
  \textbf{Rule 6:} Track section word counts from written sections\;
  \textbf{Rule 7:} On successful tool execution, recover confidence by 0.05\;
}

Sort actions by priority: \texttt{force\_finalize} $>$ \texttt{block\_publication} $>$ \texttt{switch\_profile} $>$ \texttt{inject\_directive} $>$ \texttt{continue}\;

\Return{Evaluation with prioritized actions, updated state}\;
\end{algorithm}

\subsection{Prompt Profiles}

The \governor can switch between 7 prompt profiles (\file{governor/promptProfiles.ts}, 160 lines):

\begin{table}[H]
\centering
\caption{Governor prompt profiles.}
\label{tab:governor-profiles}
\begin{tabularx}{\textwidth}{lXcc}
\toprule
\textbf{Profile} & \textbf{Trigger} & \textbf{Max Figs} & \textbf{Simple?} \\
\midrule
\texttt{default} & Initial state & 15 & No \\
\texttt{debug\_recovery} & Python errors & 8 & Yes \\
\texttt{conservative} & Data quality issues & 10 & Yes \\
\texttt{latex\_fallback} & Compilation failures & 10 & No \\
\texttt{revision\_mode} & Weak Popper review & 12 & No \\
\texttt{jf\_rewrite} & Writing phase & 15 & No \\
\texttt{academic\_figures} & Analysis phase & 15 & No \\
\bottomrule
\end{tabularx}
\end{table}

\subsection{Publication Verdict}

The \governor issues the final publication verdict based on Popper review scores:

\begin{equation}
\text{verdict} = \begin{cases}
  \text{allow} & \text{if score} \geq 80 \\
  \text{allow\_with\_note} & \text{if } 65 \leq \text{score} < 80 \\
  \text{revise} & \text{if } 45 \leq \text{score} < 65 \\
  \text{block} & \text{if score} < 45 \\
  \text{block\_fatal} & \text{if confidence} < 0.2
\end{cases}
\label{eq:pub-verdict}
\end{equation}

\subsection{Decision Logging}

All \governor decisions are recorded in a \texttt{DecisionLog} (\file{governor/decisionLog.ts}, 68 lines), an in-memory accumulator that:
\begin{itemize}
  \item Assigns each decision a UUID and timestamp.
  \item Records the pipeline stage, trigger, action, reasoning, and confidence delta.
  \item Supports serialization to/from JSON for database persistence.
  \item Tracks intervention count (excluding ``continue'' decisions).
\end{itemize}

\section{Popper --- The Peer Reviewer}
\label{sec:agents:popper}

\subsection{Role and Responsibilities}

\popper conducts AI-powered academic peer review, named after Karl Popper and his philosophy of falsificationism.
It evaluates papers across 6 dimensions and issues a publication recommendation.

\subsection{Review Process}

Implemented in \file{peerReviewService.ts} (204 lines):

\begin{enumerate}
  \item The full paper content (title, abstract, \LaTeX{} sections, BibTeX) is submitted to Claude.
  \item A detailed system prompt emulates a top-tier finance journal reviewer.
  \item Claude evaluates 6 dimensions on a 0--100 scale:
  \begin{itemize}
    \item \textbf{Originality} --- Novelty of the research question and approach.
    \item \textbf{Methodology} --- Rigor and appropriateness of methods.
    \item \textbf{Results} --- Strength and interpretation of findings.
    \item \textbf{Writing Quality} --- Clarity, organization, and prose quality.
    \item \textbf{References} --- Completeness and relevance of citations.
    \item \textbf{Reproducibility} --- Whether the paper can be replicated.
  \end{itemize}
  \item An overall score (0--100) and recommendation are produced:
  \begin{itemize}
    \item \textbf{Accept} ($\geq 80$)
    \item \textbf{Minor Revision} ($65$--$79$)
    \item \textbf{Major Revision} ($45$--$64$)
    \item \textbf{Reject} ($< 45$)
  \end{itemize}
  \item Section-by-section feedback identifies specific strengths and weaknesses.
\end{enumerate}

\subsection{Configuration}

\popper uses low temperature (0.3) for evaluation consistency and is called via the Anthropic API with the same model (\claude) as the main research pipeline.
The minimum publication threshold (default 45/100) is configurable via the admin interface.

\subsection{Failure Modes}

\begin{itemize}
  \item \textbf{Inconsistent scoring:} Despite low temperature, LLM evaluations can vary across runs.
  \item \textbf{Self-evaluation bias:} The same model generates and reviews the paper, creating potential positive bias.
  \item \textbf{Superficial review:} With limited context, \popper may miss methodological errors that a domain expert would catch.
\end{itemize}

\section{Agent Interaction Patterns}
\label{sec:agents:interaction}

\subsection{Sequential Handoff}

The primary interaction pattern is \textbf{sequential handoff}: agents execute in order (Gauss $\to$ Turing $\to$ Newton $\to$ Fisher $\to$ Euler) with each agent's output becoming the next agent's context.
Since all agents share a single Claude context window, there is no explicit message passing---the accumulated conversation history serves as the shared memory.

\subsection{Governor Intervention}

The \governor can interrupt the sequential flow by injecting directives.
For example, if \fisher encounters a Python convergence failure, the \governor may:
\begin{enumerate}
  \item Switch the profile to \texttt{debug\_recovery}.
  \item Inject a directive: ``Simplify analysis. Use basic OLS instead of GMM. Limit figures to 8.''
  \item If retries are exhausted, emit \texttt{force\_finalize} to skip remaining analysis and proceed to writing.
\end{enumerate}

\subsection{Popper Feedback Loop}

In revision mode, \popper's feedback can trigger a second pass:
\begin{enumerate}
  \item \popper reviews the paper and scores it $<65$.
  \item The \governor switches to \texttt{revision\_mode} profile.
  \item \euler re-enters writing mode to address specific criticisms.
  \item The paper is recompiled and re-reviewed.
\end{enumerate}

However, the current implementation limits revision to a single pass to prevent infinite loops.
